\section{Guenon Reader of Nietzsche}

\label{sec:GuenonReader}

\begin{quotex}
Only literature can give you that feeling of contact with another human spirit, with the totality of that spirit, his weaknesses and grandeurs, his limitations, his pettiness, his obsessions, his beliefs; with whatever moves, interests, excites, or disgusts him. Only literature can grant you access to a spirit of a dead man in a more direct, more complete, and deeper way than you would even have in conversation with a friend---as deep, as lasting as in a friendship. We never open up in a conversation as completely as when we face a blank page, addressing ourselves to an unknown reader.
\flright{\textsc{Michel Houellebecq}, \emph{Submission}}
\end{quotex}


In his novel, \emph{Submission}, \textbf{Michel Houellebecq} refers to a philosophy dissertation, held at the Catholic University of Louvain-la-Neuve, entitled \emph{Guénon lecteur de Nietzsche}\footnote{\emph{Guenon Reader of Nietzsche.}} by Robert Rediger. Over the past several months, I have made several efforts to track down a copy. A dear friend, an employee of the State Department stationed in Brussels, managed to locate a photocopy for me. When she stopped by on the way to visit her family stateside for Christmas, I was able to inspect it during her layover at the airport. Although there was not much time, I tried to extract its main themes.

The first thing to point out is the deep influence that Rediger's dissertation had on Houellebecq himself. This is shown in the main themes of the novel as expressed through the voice of Francois, his alter ego. Furthermore, the dissertation brought out several inchoate themes that were filled out in Rediger's later writings. Obviously, Rene Guenon was never a ``reader'' of Nietzsche. Rediger came to that conclusion when he wrote:

\begin{quotex}
I don't think Guenon was influenced by Nietzsche especially. His rejection of the modern world was just as vehement as Nietzsche's, but it had radically different sources.
\end{quotex}


\paragraph{Decadence}


So it is not the sources that are of interest, but rather how Guenon and Nietzsche reached similar conclusions. Specifically, there was the recognition of the decadence of Europe, the failure of Christianity to forestall it, and the rise of Islam. Francois notes this:

\begin{quotex}
The facts were plain: Europe had reached a point of such putrid decomposition that it could no longer save itself, any more than fifth-century Rome could have done. This wave of new immigrants, with their traditional culture---of natural hierarchies, the submission of women, and respect for elders---offered a historic opportunity for the moral and familial rearmament of Europe. These immigrants held out the hope of a new golden age for the old continent.
\end{quotex}


The USA is likewise in a state of decomposition, although the particular circumstances are quite different. The main difference, of course, is that the European more or less recognizes his situation while his American counterpart is blithely unaware of his real situation. Even those who call themselves conservatives in the USA are in the dark. Right wing talk shows are in a tither about Muslim violence although they, of all people, should be more concerned about those who can kill the soul rather than those who kill the body. Hence, they condemn Islam for everything in it that is traditional. For example, they emphasize the West's concern for women's and gay rights, while praising the separation of any spiritual authority over political matters.

Rediger wryly pointed out:

\begin{quotex}
Their irrational hostility to Islam should blind them to the obvious: on every question that really mattered, the nativists and the Muslims were in perfect agreement. When it came to rejecting atheism and humanism, or the necessary submission of women, or the return of patriarchy, they were fighting exactly the same fight.
\end{quotex}


\paragraph{Patriarchy}


This shows up first of all in the end of patriarchy in the West. In Francois' words, almost a direct quote from the dissertation:

\begin{quotex}
At least patriarchy existed. I mean, as a social system it was able to perpetuate itself. There were families with children, and most of them had children. In other words, it worked, where now there aren't enough children, so we're finished.
\end{quotex}


Liberalism could overturn all the established institutions: the churches, education, the workplace, political leaderships, and so on. However, the attack on the family cannot be maintained, for without children, no society can continue to exist.

Nevertheless, liberalism is adamant, while birthrates plummet. There is no longer an organic community but rather a heap of unrelated consumers. Francois lists some examples:

\begin{itemize}


\item The eco-responsible bobo

\item The show-off bourgeois woman

\item The gay-friendly nightclubbing girl

\item The satanic geek

\item The techno-Zennist
\end{itemize}


In other words, everyone is a ``type'', trying overly hard for a sui generis identity. The reader can identify more. For example, there is a facebook group called ``Liberal and proud of it''. Thus, it is not a well-thought out worldview, but rather a tribal identity marker.

Francois, despite himself, is a ``type'', perhaps the aging professor or the effete intellectual. His own life exemplifies the sterility of the modern world. Rejecting family life, he mates with his students, or with call girls when necessary. The sex scenes are not intended to be erotic, but rather sad. His preference is for anal sex or else he is ejaculating in some woman's mouth -- anywhere except where a well-bred man would consider healthy and normal.

\paragraph{Assimilation and Transformation}


Rediger then takes up a theme mentioned years ago on Gornahoor in The Prophecies of Guenon\footnote{\url{https://gornahoor.net/?p=35}}. In the \emph{Crisis of the Modern World}, Guenon wrote that only Catholicism could restore Tradition to modern world. Rediger agreed with Guenon about the medieval tradition:

\begin{quotex}
The greatness of medieval Christendom, whose artistic achievements would live forever in human memory; but little by little it had given way, it had been forced to compromise with rationalism, it had renounced its temporal powers, and so had sealed its own doom.
\end{quotex}


According to Guenon, Church leaders had forgotten the deeper truths of Tradition. Rediger goes further, asserting that the Church itself is actively engaged in subverting Tradition. He wrote:

\begin{quotex}
By dint of the affectations, the playful caresses, and the shameful fondling of the progressives, the Church had lost its ability to oppose moral decadence, to renounce homosexual marriage, abortion rights, and women in the workplace.
\end{quotex}


Even Guenon, after having failed to make any impact on Catholic thought, openly converted to Islam. Although Guenon claimed he did it for personal reasons, Rediger sees a larger perspective:

\begin{quotex}
Today this fight to establish a new organic phase of civilization could no longer be waged in the name of Christianity. Islam, its sister faith, was newer, simpler, and truer.
\end{quotex}


As Guenon predicted, the available choices are degeneration, assimilation, and transformation. Degeneration cannot continue unabated. In the novel, Houellebecq opts for a gentle assimilation and then a transformation. He foresees an alliance between the Left and the Islamist elements. That is not so far-fetched since liberalism believes it can use Islam to further erode Christian influence. However, in Houellebecq's scenario, it is Islam that assimilates the Left.

\paragraph{Polygamy and the Social Order}


Guenon was single-pointedly focused on the metaphysical elements, so he provides little or no insight into actual Muslim religious life and practice. Rediger, on the other hand, does expand on those influences, while, at the same time, adding a distinctive Western element to Islam. For this part of the dissertation, he relies more on Nietzsche than on Guenon.

Then, Rediger, also following Nietzsche, claims that Christianity is fundamentally a feminine religion. Rediger quotes Nietzsche from the
\emph{Anti-Christ}:


\begin{quotex}
If Islam despises Christianity, it has a thousandfold right to do so; Islam at least assumes that it is dealing with men.
\end{quotex}


Rediger attributes the ideals of humanism and the rights of man to the dogma of the Incarnation. Now you may object that Medieval Christianity, which Rediger lauded, was itself patriarchal and masculine. Can anyone accuse Constantine, Clovis, Charlemagne, Arthur, or Boucicaut\footnote{\url{https://gornahoor.net/?p=277}} of being feminine? Moreover, the medievals admired the great pagan warriors like Hector, Alexander, Scipio, and Julius Caesar. Even St Paul regarded effeminacy as a sin. At this point, Rediger makes the only reference to Julius Evola that I noticed: he likewise claims that this masculine element in Medieval Christianity was due solely to the still existing remnants of paganism, but not to anything specifically Christian.

First of all, Rediger accepts Nietzsche's classification of life-affirming and life-denying religions, of which Islam belongs to the former. He writes:

\begin{quotex}
Islam accepts the world, and accepts it whole. It accepts the world \emph{as such}, Nietzsche might say. For Buddhism, the world is dukkha, unsatisfactoriness, suffering. Christianity has serious reservations of its own. Isn't Satan called ``the prince of the world''? For Islam, though, the divine creation is perfect, it's an absolute masterpiece. What is the Koran, really, but one long mystical poem of praise? Of praise for the Creator, and of submission to his laws?
\end{quotex}


Rediger again turns to Nietzsche to justify the Islamic hierarchical and patriarchal order. There would be a small cadre of aristocrats on a large base of common people. Destitution would be reduced due to the requirement of alms-giving. Also, the extended family would be the ``first responders'', so to speak, in the event of the tragedies of life. These customs would reduce the size of the welfare budget. Unemployment would be reduced as more women left the workplace to stay at home, raising children.

Polygamy, too, is understood as the will to power. The stronger and the most genetically fit would have more wives, hence more children. Of course, this means that some men will have no offspring at all. This brings to mind what we recently wrote about Social Surgery\footnote{\url{https://gornahoor.net/?p=8253}}. To what extent does the community have the right to regulate the genetic makeup of its members, assuming the goal is the common good?

At this point, Rediger becomes less convincing since he does not take his ideas to their logical conclusion. He makes a good point that men of all types and castes will select the same kind of women. For example, in mate selection, most men, regardless of intelligence or station in life, would prefer a Taylor Swift. Women, on the other hand, are more malleable in their mate selection; for example, she may overlook a man's looks or age if he is wealthy or powerful.

Moreover, woman's attraction can be trained. Rediger, thus, wants to define intellectuals as ``alpha males'', deserving of the better quality women. It is beyond belief that, in the West as it now is, that women will be attracted to studious Thomists rather than the team quarterback. However, on second thought, female students are often attracted to their male professors. Asian women, and to a lesser extent Latinas, are more likely than European women to find a STEM major sexually desirable.

To accomplish this, Rediger claims that there will be women who serve as matchmakers. The details of this process are fuzzy. It would be a simple matter to determine if Islamic countries are more biologically fit than Western nations. My understanding is that the opposite is sometimes the case due to the preference for cousin marriages. But perhaps Rediger's Westernized, Nietzschean version of Islam would become the norm.

\paragraph{With a Whimper}


Overall, the social structure is based on submission, beginning with submission to God. Then there is the submission to the social order, and finally women's submission to men. Houellebecq sees the \emph{Story of O} by \textbf{Pauline Reage} as the model for submission. He sees this as occurring without much incident. This is hard to believe, although, given that the \emph{50 Shades of Grey} has sold more than 70 million copies, perhaps there is a secret desire to be dominated.

Thus, the transition occurs not with a bang, but with a whimper. Why would it be that way? Some would say it is a matter of convenience, as the intellectuals and others take advantage of the new opportunities that open up to them. Rediger concludes with a different answer:

\begin{quotex}
In the end, it was a mystery; God had ordained it so.\footnote{Some of the translations from the French text are my own.}
\end{quotex}


If that is the case, then the mystery of the three rings\footnote{\url{https://gornahoor.net/?p=7872}} will have been solved.


\begin{center}* * *\end{center}

\begin{footnotesize}
\begin{sffamily}
\texttt{aegishjalmer on 2015-12-14 at 00:40 said:}

The supposed femininity of Christianity has been asserted by some of these thinkers. I assume they refer to the historical confessions and cultures. Can anyone explain what, if anything, is specifically feminine about the essence of Christianity?

\hfill

\texttt{Max on 2015-12-14 at 10:08 said:}

On the immigrant situation

In practice, most Muslim immigrants are hardly traditional. Some of the first generation might be so to a degree, but it is the most liberal and westernized who leave their home countries in the first place. The next problem is that they do not manage to assimilate European civilization, and becomes distanced from their origin, which places them in a difficult position. This is even more pronounced in the second generation, who never received any organic influences from their own part of the world, and does not have a connection to our civilization either. It turns out that second generation ``muslims'' are if possible even more rootless and lost than the indigenous population.

Those currently fleeing to Europe does not seem to seek the restoration of any kind of ``Golden Age'' for the continent, their main motivation is rather to live an easier westernized life without the presence of war. For example, some say that the reason they fled was to avoid being drafted to the army, hardly a noble posture. Perhaps they should make an attempt to sort out their own countries instead. If the material living conditions of Europe deteriorated, they would probably not hesitate to move back.

The immigrants think that they will be able to live an organic life similar to what they had in the Middle East, but with the addition of a Western standard of living. In reality they find themselves at the bottom of the societal ladder and with a lost sense of community. Material richness is not an absoluite value, always being relative. So even if they are richer than in their old countries, they are much poorer compared to the new surroundings, and so begins to complain that is unfair and discriminating, even though they are a large cost for a society that provides them with eveything. This country produces less than 50\% of its food as it is, and uses a lot of energy simply for heating. Any leftists or ``greens'' should acknowledge that artificially increasing population in northern Europe is not very clever by their own standards of thinking, but they are not known for their coherence.

I find myself asking; in the liberals mind, who exactly is supposed to live in the Middle East? Because following their argumentation, that area of the world is made out to be completely unfit for human habitation. Therefore it could be a good idea for them to give it away for free to someone who still has use for a country of their own once it has been evacuated. In any case, a small comfort is that once enough immigrants have arrived here, either it will become so bad that no-one will want to come anymore, or the immigrants more xenophobic mentality will put an end to further immigration. A strong trend we see globally is decreasing differences between countries and continents, however not necessarily between individuals. Basically the rich parts of the world will be brought down a notch, and maybe it will further the possibility of reaching global political agreements.

If one were to have a realistic mass-immigration policy, and all of our ``empty land'' is so ideal for habitation as the politicians say, we should have some kind of settler policy like in the old American west: if you manage to carve out a living in the wilderness with your own hands, then you can have it and stay. People up here in the north respects efforts more than words; it should help integration greatly. They will also get a first-hand experience of why we are so ``distant and cold'' as they say, -``cultural enrichment'' should go both ways? While fleeing from war (in the narrative) these southerners complain about being housed in locations too remote out in the ``wilderness''. With the human races of the world adapted to specific environments, we must not forget that they are largely descended from the first urban cultures, and it would be much better for everyone involved if they could ``flee'' to an area that is as similar to theirs as possible, such as a neighboring country.

The bottom line as I see it, is that we as Europeans live in a mess of our own making, while any immigrants would live as foreigners in someone elses mess, which is arguably a much worse position. The right way to go is un-entanglement of the situation, not ``cutting the knot'', and I highly doubt that Muslims could accomplish that by themselves, which means outright replacement is not an option. The case of assimiliation means Islam would come to be something other than what we currently know it as. Islam could not conquer Europe without profound adaption, and immigrants residing in northern Europe would in some ways be at a severe disadvantage, and not fare well, in the case we abandoned the modern wellfare system for a more Islamic and/or Traditional way of life. It could just as well be what remains of Christian Europe that learns from Islam and comes out stronger.

\hfill

\texttt{Boreas on 2015-12-14 at 12:58 said:}

@ aegishjalmer: Especially one thing comes to my mind, and that is the strong separation in Christianity between creator and the creature, which leaves the personality in a feminine role in regards transcendence. I'm aware that this doesn't apply to such high mystics (initiates?) as for example Meister Eckhart, but the general tendency of Christianity is without a doubt passivity towards Being. This is one reason why I personally rejected Christianity after considering it for a while. Even the Orthodox concept of Theosis emphasizes that man and God remain separate.

\hfill

\texttt{1234kid on 2015-12-14 at 19:29 said:}

@boreas, doesnt this separation exist in islam?


\hfill

\texttt{Boreas on 2015-12-14 at 20:24 said:}

@ 1234kid: as far as I can say and know, yes it does. Sufism seems to contain some higher elements and possibilities (like does esoteric Christianity), but yes, the general tendency in all the three great monotheisms especially in their exoteric manifestations is the same.

This again reminded me of this book which I have yet to read:

\url{http://www.sophiaperennis.com/books/christianity/christianity-and-the-doctrine-of-non-dualism/}


\hfill

\texttt{Ninth Card on 2015-12-15 at 17:12 said:}

I think it's too convenient to say a tradition (christianity) is feminine and then consider everything that contradicts this statement as an exception.\\ What christianity had of pre-existing traditions was not inherited randomly but because of precise spiritual influences, as Guénon explained.\\ Part of the Egyptian, Roman, Germanic and Greek spiritual traditions survived in Christianity not as an exterior body but because they were compatible with its universal principle.\\ It happened the same thing with the influx of Daoism on Mahayana Buddhism, or Tantrism on Vajirayana Buddhism.

At the highest point of the christian tradition was Christianity feminine? Not at all.\\ Confirmation itself is a knight investiture (the bishop give a soft slap for example) and makes one a ``soldier of Christ'', the highest ascetics, many greatest christian initiates used warrior metaphors to describe their spiritual method.

Admitting that it was not feminine but only because of paganism is like implying that the weak and degenerated modern Christianity of today is the true Christianity, and the greatest Christianity of the Middle Age was the fake one.\\ So pseudo-marxist catholics of today and protestants are true christians because they are feminine, while the greatest saints were fake christians because they used warrior words? This is absurd.

I'm nobody to contest people with great traditional knowledge like Evola, but most of what he wrote about Christianity is just a post hoc rationalization of prejudices he had even before he actually knew about Tradition. (Just look at his letters of when he was young)\\ All the critics agains modern Church are valid, it's just wrong to also project them in the past and extend them to the true traditional Christianity of the past.

\hfill

\texttt{Avery on 2015-12-16 at 01:37 said:}

I don't think the nature of Christianity is strictly ``masculine'' or ``feminine'', but you can't deny that its nature is unique.

Daoism aims to harness the power of what is weak and meek for an esoteric (possibly initiated) elite. Christianity tries to harness the power of the weak and meek for the masses at large. Yet the history of Christendom is not one of submission but of reconquista, and conquista, in the name of the faith. Much later, the children of Christendom have shown a tendency to surrender their moral tenets in the interest of embracing difference and righting perceived injustices. What happened here?

I don't think Rediger or Houellebecq really plumb this idea to its depths. They are both clever men, but they can only see the psychological and demographic consequences of current visible trends.

\hfill

\texttt{Tom Walker on 2015-12-16 at 11:50 said:}

I don't think that the Knights Templar or the Christian warriors at Poitiers 732 , Malta 1551 and 1565, Lepanto 1571 , Vienna 1529 and 1583 ( to name but a few battles against Muslims ) would regard Christianity as relatively '' feminine'' . Neither would their foes . The so-called feminity of Christianity is an Evola trope with no basis in reality .\\ Anyway , all souls are passive or feminine with respect to the Spirit.

\hfill

\texttt{Alistair Fraser on 2015-12-16 at 18:34 said:}

@Max -- a very thoughtful assessment below

\hfill

\texttt{1234kid on 2015-12-16 at 20:12 said:}

@boreas, then in what way is islam addressing men where christianity is not?

i know this isn't your stance personally, im just trying to wrap my head around the idea of how a nietzschean islam would be a masculinizing force of good in the western world. two wrongs don't make a right, do they??

\hfill

\texttt{Boreas on 2015-12-21 at 10:32 said:}

@ 1234kid: well, I'm not the one defending that position.

***

As regards all souls being passive in relation to being, there is a very revealing passage in `The Doctrine of Awakening', page 86-87:

``... At this point the antitheses build up to a cosmic and titanic grandeur ending with the most paradoxical reversal of the point of view that is prevalent in Western reli\-gions. In fact, while the desire of surpassing the very Lord of creation, from this point of view, appears as something diabolical, the Buddha, instead, finds a diabolical plot in the exact opposite, that is in the attempt to stop him in the region of being, to make this region an insuperable limit, beyond which it is both absurd and mad to seek a higher liberation, Here it is the Malign One in person who urges the belief that the personal God, the God of being, is the supreme reality, and who threatens the Bud-dim with the damnation that is supposed already to have claimed other ascetics. And in another text his temptation consists of inducing the Buddha to confine himself to the path of good works, rites and sacrifices-to the path of theistic religions. But the Buddha discovers the plot, and speaks thus to M?ra: ``Well I know you, Malign One, abandon your hope: `He knows me not'; you are Mara, the Malign. Arid this Brahma here, 0 Malign One, these gods of Brahma: they are all in your hand, they are all in your power, You. O Malign One, certainly think: `He also must he in my hand, in my power!' I, however, 0 Malign One, am not in your hand, I am not in your power.''\\ There follows a symbolical test. The personal God, the Hebraic ``I am that I am,'' the God of being, whose essence is his existence, as such, cannot not be, that is, he is bound to being, he is passive with respect to being. He has not the power to go beyond being. It is here that the test occurs. Who can ``disappear'? That is, who is lord both of being and of nonbeing? Who rests neither on the one nor on the other? Brahm? cannot disappear. Instead, the Buddha disappears. All the world of Brahm? is amazed and recognizes ``the high power, the high might of the ascetic Gotama.'' Limitation is removed. The dignity of the atideva, of one who goes beyond the world of existence itself, not to mention the ``celestial'' worlds, is demonstrated. It is only left to Mara the malign to try in vain to dissuade the Buddha from spreading the doctrine\\ Here, then, is one of the extreme points in the test of the vocations: not to crave even the highest of all lives''-not only to pass from this shore to the other, but to apprehend that which lies beyond both. The words of the Awakened One are: ``Na\-ture, the gods, the lord of generation, Brahm?, the Resplendent Ones, the Powerful Ones, the Ultrapowerful Ones, all things, I have known, how unsatisfying are all things: this have I recognized and I have renounced all things, abdicated from all things. detached myself from all things, forsworn all things, disdained all things. And in this. 0 Brahma, not only am I your equal in knowledge, not only am I not less than you. but I am far greater than you,'' And the words: ``This is not mine, This am I not, this is not my self'' must be said by the ``noble son'' for the whole of that world too. It is still ``sams?ra.''\\ Is a higher limit than this conceivable? For the Ariya it is conceivable. Attach\-ment, dependence, and enjoyment are to be eradicated also in respect of the supreme goal of the Buddhist ascesis, that is to say, of extinction. Here is the final temptation and the final victory. Here the will for the unconditioned approaches the paradoxical. The ultimate truth of the series is this: he who thinks extinction, he who thinks of ex\-tinction, he who thinks ``extinction is mine'' and who rejoices in extinction-this man does not know extinction, does not know the path, is not to be counted among the ``noble disciples.'' Even in this region to feel desire and attachment-it maybe a sublimated one-means not to realize the place and signification of real liberation... ''

\url{http://www.cakravartin.com/wordpress/wp-content/uploads/2008/11/the_doctrine_of_awakening.rtf}


\hfill

\texttt{Olavsson on 2015-12-21 at 12:22 said:}

Excellent commentaries, Cologero. I've been thinking of reading the novel `Submission' for a good while. While the `ideal' scenario for the West of course is a complete restoration of sacred order based on Europe's own tradition, it is certainly true that once a civilization has so utterly and completely betrayed the Spirit itself as has been the case with the West, it should not at all be surprising if that very same Spirit is forced to reestablish itself from without, so to speak. I wouldn't hesitate for a moment to choose a Europe integrated with a complete Islamic order over a Europe abandoned to modernity. This, many `identitarians' today would no doubt see as treason. But I ask, what is worse to any noble man: betrayal of timeless divine order as such, or betrayal of some temporal collective identity? Essentially, there is only one war---the spiritual war---and two positions: for or against the integration of human life into the ascending path aligned with transcendental order.

However, it must be remembered that in practical reality the alternatives aren't that clear-cut. For example, the actual Islamic states that exist today do not represent authentic and complete traditional orders. Islamic civilization has degenerated: politically and socially; religiously, through compromises with modernity; and the `radical' Islam most loudly advocating the restoration of a traditional and imperial islamic order consists largely of violent fanatics with an inferior understanding of the tradition. As for the solution to modern disorder in Europe, there are no promising prospects in sight worthy of optimism, whether from Islam or from anywhere else. (This, of course, should cause distress to no one, as we must cultivate detachment from any and every worldly passion.) We should also keep in mind what Guénon wrote on the advent of a `Counter-Tradition' late in this cycle, which superficially would appear to be a turn away from materialistic modernity as we know it, but having the character of a grotesque parody of tradition more than anything else. (`Reign of Quantity and the Signs of the Times'). The reasons and signs supporting such a future are much more numerous today than they were in Guénon's time, even looking aside from the metaphysical consideration. In any case, those new `rightists' and identitarians who claim to oppose modernity should shut up their opportunistic whining about ``Islam'' and get themselves some proper principles, knowing what they are really `defending' other than a mixture of confused ideas and sentiments. I'd recommend checking out the self-designated reactionary and `true right' website Righton.net for some rather amusing examples.

The quote and commentary on how our contemporaries all more or less fall into various ridiculous identity-``types'' is important to take note of. This even happens to people who claim to be anti-modern and interested in Tradition, and I must admit to previously having fallen into this trap to some extent. The only way to truly avoid this danger is to practice complete detachment and cut off all personal and collective psychological bonds that obscure a pure and naked spirituality that is allowed to become the only guiding star. Easier said than done, of course.

The following reflections are pretty obvious; I only state it in case somebody less informed than Gornahoor's core community should read it:

I always thought Nietzsche's dichotomy between ``life-affirming'' and ``life-denying'' religions was nonsensical. If you are confined to a gloomy cave where you're deprived of the absolute light of metaphysical illumination and freedom, are you ``cave-denying'' if you see it as an inferior mode of being and try to forge a way out of it? The Buddha or the Christ weren't ``life-denying'' --- I would say that ``life-transcending'' is a far more appropriate way of putting it. To think that Buddha's doctrine, for example, didn't ``accept'' the world ``as such'' is only possible because Nietzsche's understanding of Reality as such is far, far more narrow and limited than the gnostic vision of such masters. Regarding those who want to remain chained to ``nature'', however, I say: let them! Nothing is easier. But they should know that it is this ``life-affirming'' attitude that is feminine, even infantile, as in a being that is too attached to ever emancipate himself from the Great Mother that is cosmic life in a pantheistic sense --- the cyclical generations of samsaric existence. The supreme expression of `masculine', heroic spirituality, on the other hand, is kaivalya, to use a yogic term, the absolute detachment, liberation, the absolute self-abiding, immanent transcendence of ultimate, luminous reality. Truthfully, I see no inherent contradiction between this and Islam's positive attitude to Allah's cosmic revelation of ``the world'', or rather, of manifestation, which in Tradition embraces many more existential states of being than the immediate sensible reality we take for ``the world.'' Even in the generic Hindu tradition, whose universal aim is moksha, supreme deliverance from conditioned existence, there is still traditionally a coexisting `affirmation' of life that nearly surpasses that of Islam itself. Both masculine life-transcending and feminine life-affirming religious values can in fact coexist in the traditional order, all values in their proper place. And so, even in Islam, the impetus to spiritual transcendence has been strongly present. After all --- and Rediger seems to contradict himself partly here --- the emphasis of Islam is upon Allah, the Absolute, not upon ``the world''; and while Allah is present within ``the world'', He is also radically transcendent to it. This gives an inferior value to ``the world as such'' compared to the uncreated Reality of Allah.

@ Boreas: What spiritual path are you pursuing now if you rejected Christian Orthodoxy? Have you decided on a specific course of action?

Whether or not the ``man and God remaining separate'' (even in Theosis) must be problematic, depends on how it's understood. At least it isn't more problematic than the neo-Vedantist ideas critiqued somewhere on here about the highest transcendence being a sort of ``merging'' with a cosmic ``Oneness'', disappearing in an impersonal sea. Certainly the overcoming of all dualism is required to participate in the most absolute `state', but perhaps absolute unity (with Godhead) and absolute `Purushic' separation or isolation of the Self is not as contradictory as some may think. Reminds me of the Vajrayana sage and initiate Dolpopa who said somewhere that ultimate reality is neither One nor many, being limited by neither. Such discussions often become overly abstract and speculative, much because what we're trying to understand in verbal terms here essentially belongs to a reality that is inexpressible (the initiatic `secret') and must be `experienced' directly, as real transcendence, for infallible knowledge to be possible. This exactly was the reason for the Buddha Shakyamuni's infamous `silence' on such matters. But as a general rule, if a tradition excludes the possibility of `realizing' a reality that is unborn, uncreated, unconditioned and so on, then it doesn't go all the way to the peak of the `Greater Mysteries' (see Guénon on the distinction between the `Lesser' and `Greater' Mysteries).

Regarding the question on the femininity and passivity of the soul in respect to Being or Spirit: This is indeed true for the soul as such. The point in Buddha's asceticism, though (since Boreas quoted a passage from Evola's `Doctrine of Awakening') is to go beyond the soul as such since it is conditioned. What is inherently conditioned cannot ``become'' unconditioned -- and if it could, it would no longer be the same thing. This is true of the soul. One can, however, ``awaken'' to the Tathata, Dharmakaya, the absolute nature that has always been unchanging and ``behind'' the obscurations and appearances of conditioned manifestation. But this does not mean that a *soul* strictly speaking suddenly becomes superior to the dimension of Being. Overcoming the limitations of the soul is possible because the soul is not our true, unborn nature that was before the world, before everything.

\hfill

\texttt{Boreas on 2015-12-21 at 15:20 said:}

@ Olavsson: Good points regarding the soul and being. I stand corrected.

As to what my path is after rejecting Orthodoxy, I'm still undecided. Currently I'm trying to become an enlightened Pagan and I'm studying different heathen traditions. It might be that I'm arriving at some sort of ``esoteric synthesis'' (not syncretism) outside any established tradition. Maybe a term Christo-Pagan could be a fitting one for me personally. I'm avoiding the neo-pagan trap however, and would not hesitate to join Christian forces to establish ``a social reign of Christ'' (I'm still essentially a Monarchist), but that seems to be utopian thinking at the present moment, since all that the orthodoxes of my native land for example are talking socially is democracy, multi-culturalism and inter-faith relationships. I'm certainly not an anti-Christian though, since I value the teachings of Christ very highly.

\hfill

\texttt{ARYANWOLF5359021 on 2016-05-13 at 23:15 said:}

``Nietzsche, with his prophetic awareness -- in an intellectual sense -- foresaw this coming disaster ... Man has been downgraded. He is now already `sub'. We have been made sub-human. So we must take the Nietzschean image of reaching the Übermensch as an Islamic duty, an Islamic call.'' --- Shaykh Abdalqadir as- Sufi

\url{https://en.wikipedia.org/wiki/Abdalqadir_as-Sufi}


\hfill

\texttt{Cologero on 2016-05-14 at 09:49 said:}

Time will tell, Aryan Wolf.

\end{sffamily}
\end{footnotesize}

